\documentclass[11pt]{article}
\usepackage{amsmath,amssymb,booktabs,geometry}
\geometry{margin=1in}

\title{Two Coherent Vacuum Plane Waves at $120^\circ$:\\
Vector Superposition and the Atom 2.1 Interpretation Boundary}
\author{[Authors to be supplied]\\[Affiliations to be supplied]}
\date{Working manuscript v0.1.0 --- 2026-09-01}

\begin{document}
\maketitle

\begin{abstract}
We derive the classical vector superposition of two mutually coherent,
equal-frequency, equal-amplitude vacuum plane waves crossing at $120^\circ$.
For common polarization normal to the wavevector plane, exact field, energy,
Poynting-flux, momentum-density, and invariant relations are obtained as functions
of local relative phase. A phase of $120^\circ$ is not an extremum. Linear Maxwell
theory predicts interference and no new photon or neutrino. The Atom 2.1 term
``neutrino seed'' is therefore retained only as an unvalidated hypothesis label.
\end{abstract}

\noindent\textbf{Keywords:} electromagnetic interference; Maxwell equations;
plane waves; 120-degree geometry; phase control; hypothesis boundary

\section{Introduction and Research Question}
What field, energy, and momentum pattern follows from linear Maxwell theory when
two identical coherent vacuum modes cross at $120^\circ$, and is a local phase of
$120^\circ$ distinguished? ``Photon plane wave'' is used informally for a classical
mode; QED and particle creation are outside scope.

\section{Background}
Maxwell's vacuum equations are linear, and coherent field sums exhibit standard
interference~\cite{jackson1999,bornwolf1999,mandelwolf1995}. This established result
is separated from the Atom 2.1 proposal to interpret a selected region as a
``neutrino seed.''

\section{Methods and Reproducibility}
Let $k=\omega/c$ and choose
\begin{align}
 \mathbf{n}_1&=(1/2,\sqrt{3}/2,0), &
 \mathbf{n}_2&=(1/2,-\sqrt{3}/2,0), &
 \mathbf{k}_i&=k\mathbf{n}_i. \tag{1}
\end{align}
Then $\mathbf{n}_1\cdot\mathbf{n}_2=-1/2$, so the crossing angle is
$2\pi/3$, and
\begin{equation}
 \mathbf{k}_1+\mathbf{k}_2=k\hat{\mathbf{x}},\qquad
 \mathbf{k}_1-\mathbf{k}_2=\sqrt{3}k\hat{\mathbf{y}}. \tag{2}
\end{equation}
For common electric polarization $\hat{\mathbf{z}}$ and independent phase offset
$\delta$, use peak phasors
\begin{align}
 \widetilde{\mathbf E}&=E_0\hat{\mathbf z}
 \left[e^{i\mathbf{k}_1\cdot\mathbf r}+e^{i(\mathbf{k}_2\cdot\mathbf r+\delta)}\right], \tag{3}\\
 \widetilde{\mathbf B}&=\frac{E_0}{c}\left[(\mathbf n_1\times\hat{\mathbf z})e^{i\mathbf{k}_1\cdot\mathbf r}
 +(\mathbf n_2\times\hat{\mathbf z})e^{i(\mathbf{k}_2\cdot\mathbf r+\delta)}\right]. \tag{4}
\end{align}
The local phase is $\Gamma=\delta-\sqrt{3}ky$, giving fringe spacing
$\lambda/\sqrt{3}$. Infinite monochromatic coherent waves, linear vacuum response,
and local cycle averages are assumed.

Reproducibility commands are
\begin{verbatim}
python nodes/N015_trinity_phase_lock_seed_event/code/sim.py
python scripts/validate_papers.py
python -m pytest -q
\end{verbatim}
No data fitting, randomness, or external dataset is used.

\section{Results}
The phasor magnitudes and cycle-averaged observables are
\begin{align}
 |\widetilde{\mathbf E}|^2&=2E_0^2(1+\cos\Gamma), &
 c^2|\widetilde{\mathbf B}|^2&=E_0^2(2-\cos\Gamma), \tag{5}\\
 \langle u\rangle&=\varepsilon_0E_0^2(1+\tfrac14\cos\Gamma), &
 \langle\mathbf S\rangle&=I_0\hat{\mathbf x}(1+\cos\Gamma), \tag{6}\\
 \langle\mathbf g\rangle&=\langle\mathbf S\rangle/c^2, &
 \langle E^2-c^2B^2\rangle&=\tfrac32E_0^2\cos\Gamma, \tag{7}
\end{align}
where $I_0=\varepsilon_0cE_0^2/2$ and
$\langle\mathbf E\cdot\mathbf B\rangle=0$.

\begin{table}[h]
\centering
\caption{Exact phase values for the $120^\circ$ crossing geometry.}
\begin{tabular}{ccccc}
\toprule
$\Gamma$ & $|\widetilde{\mathbf E}|/E_0$ &
$\langle u\rangle/(\varepsilon_0E_0^2)$ & $\langle S_x\rangle/I_0$ &
$\langle E^2-c^2B^2\rangle/E_0^2$\\
\midrule
$0$ & $2$ & $5/4$ & $2$ & $3/2$\\
$2\pi/3$ & $1$ & $7/8$ & $1/2$ & $-3/4$\\
$\pi$ & $0$ & $3/4$ & $0$ & $-3/2$\\
\bottomrule
\end{tabular}
\end{table}

The phase derivatives are proportional to $-\sin\Gamma$, so $2\pi/3$ is not
stationary. Averaging one fringe removes $\cos\Gamma$ and recovers the sums of the
two input energy densities and vector Poynting fluxes.

\section{Discussion}
The field retains both wavevectors; Eq.~(2) is geometry, not conversion to a new
mode. Interference redistributes energy and momentum locally without creating net
energy. Maxwell theory supplies no localization, stability, neutrino quantum
numbers, or independent photon. A future Atom 2.1 seed model must add explicit
dynamics and a discriminating, falsifiable observable.

\section{Limitations and Scope}
The calculation excludes finite beams, wave packets, apertures, sources, matter,
nonlinear optics, gravity, CMB forcing, QED, neutrino phenomenology, and charge-only
or magnetic-only propagating-wave claims. Only a common linear polarization is
evaluated. Internal checks are not empirical validation.

\section{Conclusion}
Two equal coherent Maxwell modes crossing at $120^\circ$ yield the standard
interference relations above. Local phase $120^\circ$ gives intermediate values,
not an extremum, and does not derive a particle.

\section*{Data and Code Availability}
No empirical data were used. Code and tests are stored with node N015 and in the
repository test suite.

\section*{Author Contributions}
[To be supplied by the actual authors; no roles inferred.]

\section*{Acknowledgments}
[None declared in this repository version.]

\section*{Declarations}
Funding and conflicts of interest are not declared and must be supplied by the
authors. Ethics approval and consent are not applicable to this theoretical work.
No empirical validation is claimed.

\section*{Version and Provenance}
Working manuscript v0.1.0, 2026-09-01. Derived from nodes N000, N011, and N015;
not peer reviewed.

\bibliographystyle{plain}
\bibliography{refs}
\end{document}
